\documentclass[12pt,a4paper]{article}

% Codificacao e idioma
\usepackage[utf8]{inputenc}
\usepackage[T1]{fontenc}
\usepackage[brazil]{babel}

% Matematica
\usepackage{amsmath, amssymb, amsthm}
\usepackage{mathtools}

% Layout
\usepackage[margin=2.5cm]{geometry}
\usepackage{setspace}
\usepackage{parskip}

% Tabelas e figuras
\usepackage{booktabs}
\usepackage{array}
\usepackage{multirow}
\usepackage{float}

% Algoritmos
\usepackage[ruled,vlined,portuguese]{algorithm2e}

% Codigo fonte
\usepackage{listings}
\usepackage{xcolor}

% Links
\usepackage{hyperref}
\hypersetup{colorlinks=true, linkcolor=blue, citecolor=blue}

% Referencias
\usepackage[numbers,sort]{natbib}

% ---------------------------------------------------------------------------
% Ambientes de teorema
% ---------------------------------------------------------------------------
\newtheorem{theorem}{Teorema}[section]
\newtheorem{lemma}[theorem]{Lema}
\newtheorem{corollary}[theorem]{Corolário}
\newtheorem{proposition}[theorem]{Proposição}
\theoremstyle{definition}
\newtheorem{definition}[theorem]{Definição}
\newtheorem{remark}[theorem]{Observação}

% ---------------------------------------------------------------------------
% Comandos matematicos
% ---------------------------------------------------------------------------
\newcommand{\bigo}[1]{\mathcal{O}\!\left(#1\right)}
\newcommand{\bigtheta}[1]{\Theta\!\left(#1\right)}
\newcommand{\bigomega}[1]{\Omega\!\left(#1\right)}
\newcommand{\C}[2]{\binom{#1}{#2}}
\newcommand{\Sk}{\mathcal{S}_{15}}
\newcommand{\Sp}{\mathcal{S}_{p}}
\newcommand{\SB}{\mathrm{SB}}
\newcommand{\OPT}{\mathrm{OPT}}
\newcommand{\Harm}[1]{H_{#1}}

% ---------------------------------------------------------------------------
\begin{document}
% ---------------------------------------------------------------------------

\begin{titlepage}
  \centering
  \vspace*{2cm}
  {\Large \textbf{PONTIFÍCIA UNIVERSIDADE CATÓLICA DO PARANÁ}\\[0.3em]
   \large Escola Politécnica --- Bacharelado em Ciência da Computação}
  \vspace{3cm}

  {\Huge \textbf{Trabalho Avaliativo --- RA03}}\\[1em]
  {\LARGE Cobertura de Combinações: Complexidade, Algoritmos e Análise}

  \vspace{3cm}
  {\large Disciplina: Complexidade de Algoritmos\\[0.5em]
   Professor: Edson Emílio Scalabrin}

  \vfill
  {\large Curitiba, junho de 2026}
\end{titlepage}

% ---------------------------------------------------------------------------
\tableofcontents
\newpage
% ---------------------------------------------------------------------------

\section{Definição Formal do Problema}

\subsection{Universo e conjuntos de combinações}

Seja o universo $U = \{1, 2, \ldots, 25\}$ com $|U| = 25$.
Definimos o conjunto de todas as combinações de tamanho $p$ extraídas de $U$:
\[
  \mathcal{S}_p \;=\; \bigl\{\, X \subseteq U \;\big|\; |X| = p \,\bigr\}
\]
cuja cardinalidade é dada pelo coeficiente binomial $|\mathcal{S}_p| = \C{25}{p}$.

\begin{table}[H]
\centering
\caption{Cardinalidades dos conjuntos do problema.}
\begin{tabular}{ccr}
\toprule
Conjunto & Fórmula & Cardinalidade \\
\midrule
$\mathcal{S}_{15}$ & $\C{25}{15}$ & 3.268.760 \\
$\mathcal{S}_{14}$ & $\C{25}{14}$ & 4.457.400 \\
$\mathcal{S}_{13}$ & $\C{25}{13}$ & 5.200.300 \\
$\mathcal{S}_{12}$ & $\C{25}{12}$ & 5.200.300 \\
$\mathcal{S}_{11}$ & $\C{25}{11}$ & 4.457.400 \\
\bottomrule
\end{tabular}
\end{table}

\subsection{Problema de cobertura (Programas 2--5)}

Para cada $p \in \{14, 13, 12, 11\}$, o problema consiste em determinar o
\emph{menor subconjunto} $\SB_{15,p} \subseteq \mathcal{S}_{15}$ tal que:
\[
  \forall\, Y \in \mathcal{S}_p,\quad \exists\, X \in \SB_{15,p} \text{ tal que } Y \subseteq X
\]

\begin{definition}[Minimum Set Cover]
  Dado um universo finito $\mathcal{U}$, uma família de conjuntos
  $\mathcal{F} = \{F_1, \ldots, F_m\} \subseteq 2^{\mathcal{U}}$ e um
  inteiro $k$, o problema \textsc{Min-Set-Cover} pergunta se existe
  $\mathcal{C} \subseteq \mathcal{F}$ com $|\mathcal{C}| \leq k$ tal que
  $\bigcup_{F \in \mathcal{C}} F = \mathcal{U}$.
\end{definition}

Nossa instância é:
$\mathcal{U} = \mathcal{S}_p$,\quad
$\mathcal{F} = \bigl\{\, \{Y \in \mathcal{S}_p \mid Y \subseteq X\} \;\big|\; X \in \mathcal{S}_{15} \,\bigr\}$.

% ---------------------------------------------------------------------------
\section{Limites Inferiores (Lower Bounds)}
% ---------------------------------------------------------------------------

\subsection{Limite por relaxação linear (LP)}

\begin{theorem}[Lower Bound LP]\label{thm:lb-lp}
  $|\SB_{15,p}| \;\geq\; \left\lceil \dfrac{\C{25}{p}}{\C{15}{p}} \right\rceil$
\end{theorem}

\begin{proof}
  Cada $X \in \mathcal{S}_{15}$ contém exatamente $\C{15}{p}$ subconjuntos de
  tamanho $p$ (os $p$-subconjuntos de seus 15 elementos).
  Como todo $Y \in \mathcal{S}_p$ deve ser coberto por pelo menos um $X \in \SB$:
  \[
    |\SB| \cdot \C{15}{p} \;\geq\; |\mathcal{S}_p| = \C{25}{p}
    \;\implies\;
    |\SB| \;\geq\; \left\lceil \frac{\C{25}{p}}{\C{15}{p}} \right\rceil. \qedhere
  \]
\end{proof}

\begin{table}[H]
\centering
\caption{Lower bounds LP para cada valor de $p$.}
\begin{tabular}{crrrrrr}
\toprule
$p$ & $\C{25}{p}$ & $\C{15}{p}$ & $\lceil \C{25}{p}/\C{15}{p} \rceil$ & --- & --- & --- \\
\midrule
14 & 297.160 & 297.172 & 532.555 & 1.79$	imes$ & --- & 16.3\% \
13 & 49.527 & 58.887 & 	extit{(em execução)} & --- & --- & --- \
12 & 11.430 & 13.175 & 38.100 & 3.33$	imes$ & --- & 1.2\% \
11 & 3.266 & 3.370 & 	extit{(em execução)} & --- & --- & --- \
\bottomrule
\end{tabular}
\end{table}

\subsection{Limite de Schönheim}

\begin{definition}[Número de cobertura]
  $C(n, k, t)$ é o menor inteiro $m$ tal que existem $m$ subconjuntos
  de tamanho $k$ de $\{1,\ldots,n\}$ cobrindo todo $t$-subconjunto.
\end{definition}

\begin{theorem}[Schönheim, 1964]\label{thm:schoenheim}
  $C(n, k, t) \;\geq\; L(n, k, t)$, onde:
  \[
    L(n, k, t) = \left\lceil \frac{n}{k} \cdot L(n-1,\, k-1,\, t-1) \right\rceil,
    \qquad L(t,\, t,\, t) = L(k,\, k,\, t) = 1.
  \]
\end{theorem}

\begin{proof}
  Por indução em $n - t$.
  Fixe $x \in \{1,\ldots,n\}$ e seja $\SB$ uma cobertura ótima com $|\SB| = C(n,k,t)$.
  Defina $\SB_x = \{X \in \SB \mid x \in X\}$.

  Os conjuntos de $\SB_x$ removidos de $x$ formam uma cobertura de $(t-1)$-subconjuntos
  de $\{1,\ldots,n\} \setminus \{x\}$, logo $|\SB_x| \geq C(n-1, k-1, t-1)$ (h.i.).

  Por argumento de média sobre $x \in \{1,\ldots,n\}$:
  \[
    \sum_{x=1}^{n} |\SB_x| = k \cdot |\SB|
    \;\implies\;
    \frac{k \cdot |\SB|}{n} = \overline{|\SB_x|} \;\geq\; L(n-1, k-1, t-1),
  \]
  pois o mínimo é atingido por algum $x$. Logo
  $|\SB| \geq \lceil (n/k) \cdot L(n-1, k-1, t-1) \rceil = L(n, k, t)$.
\end{proof}

\begin{table}[H]
\centering
\caption{Comparação entre lower bounds para $n=25$, $k=15$.}
\begin{tabular}{crrrr}
\toprule
$p$ & $|\mathcal{S}_p|$ & LB-LP & LB-Schönheim & Razão Sch/LP \\
\midrule
14 & 297.160 & 297.172 & 532.555 & 1.79$	imes$ & --- & 16.3\% \
13 & 49.527 & 58.887 & 	extit{(em execução)} & --- & --- & --- \
12 & 11.430 & 13.175 & 38.100 & 3.33$	imes$ & --- & 1.2\% \
11 & 3.266 & 3.370 & 	extit{(em execução)} & --- & --- & --- \
\bottomrule
\end{tabular}
\end{table}

O bound de Schönheim é mais apertado que o LP para $p \in \{12, 13\}$.

% ---------------------------------------------------------------------------
\section{NP-dificuldade do Problema}
% ---------------------------------------------------------------------------

\begin{theorem}[Karp, 1972]\label{thm:np}
  \textsc{Min-Set-Cover} é NP-difícil.
\end{theorem}

\begin{proof}
  Redução polinomial de \textsc{Vertex Cover} (NP-completo).
  Dado grafo $G = (V, E)$:
  \begin{itemize}
    \item Universo: $\mathcal{U} = E$
    \item Para cada $v \in V$: conjunto $S_v = \{e \in E \mid v \in e\}$
  \end{itemize}
  $C \subseteq V$ é vertex cover $\;\Longleftrightarrow\;$
  $\{S_v \mid v \in C\}$ é set cover:
  toda aresta $e = \{u, v\}$ tem $u \in C$ ou $v \in C$ (VC) $\Leftrightarrow$
  $e \in S_u$ ou $e \in S_v$ (SC). A redução é em tempo $O(|V| + |E|)$. \qedhere
\end{proof}

\begin{corollary}
  Os Programas 2--5 são instâncias de \textsc{Min-Set-Cover}, portanto NP-difíceis.
  Não existe algoritmo polinomial exato a menos que $\mathrm{P} = \mathrm{NP}$.
\end{corollary}

% ---------------------------------------------------------------------------
\section{Algoritmos Implementados}
% ---------------------------------------------------------------------------

\subsection{Programa 1 --- Geração de combinações}

\begin{algorithm}[H]
\SetAlgoLined
\KwIn{$n = 25$, $p \in \{15, 14, 13, 12, 11\}$}
\KwOut{Array $A$ de bitmasks uint32 com todos os $X \in \mathcal{S}_p$}
\BlankLine
$A \leftarrow []$\;
\ForEach{$(e_1, \ldots, e_p) \in \binom{\{1,\ldots,n\}}{p}$ em ordem lexicográfica}{
  $\text{mask} \leftarrow \sum_{i=1}^{p} 2^{e_i - 1}$\;
  $A.\text{append}(\text{mask})$\;
}
\Return $A$\;
\caption{Geração por bitmask}
\end{algorithm}

\paragraph{Representação por bitmask.}
Cada subconjunto $X \subseteq U$ é representado como inteiro de 25 bits:
elemento $e \in X$ ocupa o bit $e-1$. A verificação $Y \subseteq X$ torna-se:
\[
  Y \subseteq X \;\Longleftrightarrow\; (X \,\&\, Y) = Y
\]
operação $O(1)$ em hardware, $\bigtheta{1}$ por par — frente a $\bigtheta{p}$ com frozensets.

\subsubsection*{Análise de complexidade --- Programa 1}

\begin{itemize}
  \item \textbf{Tempo:} $\bigtheta{\C{n}{p}}$ — cada combinação é visitada exatamente uma vez.
  \item \textbf{Espaço (lazy):} $\bigtheta{p}$ por iteração; $\bigtheta{\C{n}{p}}$ se armazenado em array.
  \item \textbf{Gargalo:} geração de $\C{25}{15} = 3.268.760$ combinações com
        \texttt{itertools.combinations} e conversão para \texttt{uint32} consome $\approx 10$s.
\end{itemize}

% ---------------------------------------------------------------------------
\subsection{Solver Greedy (Programas 2--5)}

\begin{algorithm}[H]
\SetAlgoLined
\KwIn{$\mathcal{S}_{15}$, $\mathcal{S}_p$}
\KwOut{$\SB \subseteq \mathcal{S}_{15}$ cobrindo $\mathcal{S}_p$}
\BlankLine
Construir índices: $\texttt{idx}_{15}[\text{mask}] \to i$;\quad
                   $\texttt{idx}_p[\text{mask}] \to j$\;
$\texttt{count}[i] \leftarrow \C{15}{p}$ para todo $i$\;
$\texttt{nao\_coberto}[j] \leftarrow \mathbf{true}$ para todo $j$\;
$\SB \leftarrow \emptyset$;\quad inicializar heap $H$ com $(-\C{15}{p},\, i)$ para todo $i$\;
\BlankLine
\While{$\exists\, j : \texttt{nao\_coberto}[j]$}{
  $i^* \leftarrow \arg\max_i \texttt{count}[i]$ \tcp*{via heap ou argmax NumPy}
  $X^* \leftarrow \mathcal{S}_{15}[i^*]$;\quad $\SB \leftarrow \SB \cup \{X^*\}$;\quad $\texttt{count}[i^*] \leftarrow -1$\;
  \ForEach{$Y \in \binom{X^*}{p}$}{
    \lIf{$\texttt{nao\_coberto}[\texttt{idx}_p[Y]]$}{
      $\texttt{nao\_coberto}[\texttt{idx}_p[Y]] \leftarrow \mathbf{false}$\;
      \ForEach{$Z \in \mathcal{S}_{15}$ com $Y \subseteq Z$}{
        $\texttt{count}[\texttt{idx}_{15}[Z]] \mathrel{-}= 1$
      }
    }
  }
}
\Return $\SB$\;
\caption{Greedy Set Cover com heap lazy}
\end{algorithm}

\subsubsection*{Análise de complexidade --- Greedy}

Sejam:
$N_{15} = |\mathcal{S}_{15}| = \C{25}{15}$,\;
$N_p = |\mathcal{S}_p| = \C{25}{p}$,\;
$K = |\SB_{\text{greedy}}|$.

\paragraph{Pré-processamento.}
\begin{itemize}
  \item Construção dos índices hash: $\bigtheta{N_{15} + N_p}$ tempo e espaço.
  \item Inicialização do heap: $\bigtheta{N_{15}}$ (heapify em array uniforme).
\end{itemize}

\paragraph{Custo por iteração.}
Cada iteração processa $X^* \in \mathcal{S}_{15}$:
\begin{enumerate}
  \item Selecionar $X^*$: $\bigo{\log N_{15}}$ (heap) ou $\bigo{N_{15}}$ (argmax).
  \item Gerar $\C{15}{p}$ subconjuntos $Y$ de $X^*$.
  \item Para cada $Y$ recém-coberto: gerar $\C{25-p}{15-p}$ extensões de $Y$ até tamanho 15
        e decrementar \texttt{count}. Total de atualizações:
        \[
          \C{15}{p} \cdot \C{25-p}{15-p} \text{ por iteração.}
        \]
\end{enumerate}

\paragraph{Complexidade total.}
\[
  \text{Tempo} = \bigtheta{N_{15} + N_p}
    + K \cdot \left[\C{15}{p}\cdot\C{25-p}{15-p}
    + \begin{cases} \bigo{\log N_{15}} & \text{heap} \\ \bigo{N_{15}} & \text{argmax} \end{cases}
    \right]
\]

Como $\C{15}{p}\cdot\C{25-p}{15-p} \gg \log N_{15}$ para os valores considerados:
\[
  \boxed{\text{Tempo} = \bigo{K \cdot \C{15}{p} \cdot \C{25-p}{15-p}}}
\]
\[
  \text{Espaço} = \bigtheta{N_{15} + N_p}
\]

\begin{table}[H]
\centering
\caption{Parâmetros de complexidade do Greedy por iteração.}
\begin{tabular}{crrrrr}
\toprule
$p$ & $\C{15}{p}$ & $\C{25-p}{15-p}$ & Atualizações/iter & $K$ (greedy) & Estratégia \\
\midrule
14 & 297.160 & 297.172 & 532.555 & 1.79$	imes$ & --- & 16.3\% \
13 & 49.527 & 58.887 & 	extit{(em execução)} & --- & --- & --- \
12 & 11.430 & 13.175 & 38.100 & 3.33$	imes$ & --- & 1.2\% \
11 & 3.266 & 3.370 & 	extit{(em execução)} & --- & --- & --- \
\bottomrule
\end{tabular}
\end{table}

\paragraph{Seleção automática de estratégia.}
O limiar de 10.000 atualizações por iteração determina a escolha:
\begin{itemize}
  \item $p \in \{14, 13\}$ ($\leq 6.930$ atualizações): heap lazy preferível —
        as poucas atualizações custam $\bigo{\log N_{15}}$ cada, mas as muitas iterações $K$
        tornam o heap $\bigo{K \log N_{15}}$ competitivo.
  \item $p \in \{12, 11\}$ ($\geq 130.130$ atualizações): argmax NumPy preferível —
        evita os 130K+ \texttt{heappush} por iteração, que custam mais que uma varredura
        $\bigo{N_{15}}$ em NumPy.
\end{itemize}

\subsubsection*{Prova de corretude do Greedy}

\begin{theorem}[Corretude]\label{thm:greedy-correctness}
  O algoritmo Greedy termina e retorna $\SB$ tal que todo $Y \in \mathcal{S}_p$
  está coberto por algum $X \in \SB$.
\end{theorem}
\begin{proof}
  \textbf{Terminação:} A cada iteração, $|\texttt{nao\_coberto}|$ decresce
  estritamente (o $X^*$ escolhido cobre ao menos um $Y$ não coberto).
  Como $|\mathcal{S}_p|$ é finito, o loop termina em $\leq |\mathcal{S}_p|$ iterações.

  \textbf{Corretude:} Ao término, $\texttt{nao\_coberto}[j] = \mathbf{false}$
  para todo $j$, logo todo $Y \in \mathcal{S}_p$ foi coberto em alguma iteração
  por um $X^* \in \SB$. \qedhere
\end{proof}

\subsubsection*{Razão de aproximação}

\begin{theorem}[Johnson 1974; Lovász 1975]\label{thm:approx}
  Seja $\Harm{m} = \sum_{i=1}^{m} \frac{1}{i} \leq \ln m + 1$ o $m$-ésimo
  número harmônico. Então:
  \[
    |\SB_{\text{greedy}}| \;\leq\; \Harm{|\mathcal{S}_p|} \cdot |\OPT|
    \;\leq\; (\ln|\mathcal{S}_p| + 1) \cdot |\OPT|.
  \]
\end{theorem}
\begin{proof}
  Numere os elementos de $\mathcal{S}_p$ na ordem em que são cobertos pelo greedy:
  $y_1, \ldots, y_m$ ($m = |\mathcal{S}_p|$). Atribua custo
  $c(y_i) = 1/d_t$ a cada $y_i$ coberto na iteração $t$, onde $d_t$ é o número
  de novos elementos cobertos nessa iteração.

  \textbf{Passo 1 (custo total):}
  $\sum_i c(y_i) = \sum_t d_t \cdot \frac{1}{d_t} = |\SB_{\text{greedy}}|$.

  \textbf{Passo 2 (custo por conjunto ótimo):}
  Para $X \in \OPT$ com $s = |X \cap \mathcal{S}_p|$, os elementos de $X$
  cobertos em ordem $z_1, \ldots, z_s$ têm:
  \[
    c(z_j) \leq \frac{1}{s - j + 1}
  \]
  pois quando $z_j$ é coberto, ainda restam $\geq s - j + 1$ elementos de
  $X \cap \mathcal{S}_p$ não cobertos, e o greedy escolheu $X^*$ com cobertura
  $\geq s - j + 1$. Logo a contribuição de $X$ ao custo é
  $\leq \Harm{s} \leq \Harm{m}$.

  \textbf{Conclusão:}
  $|\SB_{\text{greedy}}| = \sum_i c(y_i) \leq |\OPT| \cdot \Harm{m}$. \qedhere
\end{proof}

\begin{remark}
  A razão $\Harm{m}$ é assintoticamente ótima: não existe algoritmo polinomial
  com razão $(1-\varepsilon)\ln m$ para nenhum $\varepsilon > 0$, a menos que
  $\mathrm{P} = \mathrm{NP}$ \cite{feige1998}.
\end{remark}

% ---------------------------------------------------------------------------
\subsection{Solver por Programação Inteira Linear (ILP)}

\paragraph{Modelagem formal.}
\[
  \min \sum_{i=1}^{N_{15}} x_i
\]
\[
  \text{sujeito a} \quad
  \sum_{\substack{i \,:\, Y \subseteq \mathcal{S}_{15}[i]}} x_i \;\geq\; 1
  \quad \forall\, Y \in \mathcal{S}_p
\]
\[
  x_i \in \{0, 1\} \quad \forall\, i \in \{1, \ldots, N_{15}\}
\]

\paragraph{Inviabilidade em escala completa.}

\begin{table}[H]
\centering
\caption{Dimensão do ILP em escala completa.}
\begin{tabular}{crrr}
\toprule
$p$ & Variáveis & Restrições & Coef. não-nulos \\
\midrule
14 & 297.160 & 297.172 & 532.555 & 1.79$	imes$ & --- & 16.3\% \
13 & 49.527 & 58.887 & 	extit{(em execução)} & --- & --- & --- \
12 & 11.430 & 13.175 & 38.100 & 3.33$	imes$ & --- & 1.2\% \
11 & 3.266 & 3.370 & 	extit{(em execução)} & --- & --- & --- \
\bottomrule
\end{tabular}
\end{table}

O ILP direto é inviável em escala completa. A estratégia adotada é:
\emph{usar a solução greedy como universo de candidatos e resolver o ILP
sobre esse conjunto reduzido}, encontrando o subconjunto mínimo do greedy
que ainda cobre $\mathcal{S}_p$.

\subsubsection*{Análise de complexidade --- ILP}

\begin{itemize}
  \item \textbf{Complexidade teórica:} NP-difícil no pior caso
        (pois \textsc{Min-Set-Cover} é NP-difícil pelo Teorema~\ref{thm:np}).
  \item \textbf{Tempo branch-and-bound:} exponencial no pior caso;
        na prática, depende do gap LP relaxado.
  \item \textbf{Espaço:} $\bigo{|\text{candidatos}| + |\mathcal{S}_p|}$ para o modelo.
  \item \textbf{Gargalo:} construção do mapa de cobertura em Python puro:
        $\bigo{|\mathcal{S}_p| \cdot \C{25-p}{15-p}}$ operações.
\end{itemize}

% ---------------------------------------------------------------------------
\subsection{Solver Randômico (Monte Carlo)}

\begin{algorithm}[H]
\SetAlgoLined
\KwIn{$\mathcal{S}_{15}$, $\mathcal{S}_p$, $T$ tentativas}
\KwOut{Melhor $\SB$ encontrado}
\BlankLine
$\SB^* \leftarrow \mathcal{S}_{15}$\;
\For{$t = 1$ \KwTo $T$}{
  $\pi \leftarrow$ permutação aleatória de $\mathcal{S}_{15}$\;
  $\SB \leftarrow \emptyset$;\quad $N \leftarrow \mathcal{S}_p$\;
  \ForEach{$X \in \pi$}{
    \lIf{$N = \emptyset$}{\textbf{break}}
    \lIf{$\exists\, Y \in N : Y \subseteq X$}{$\SB \leftarrow \SB \cup \{X\}$;\; $N \leftarrow N \setminus \{Y : Y \subseteq X\}$}
  }
  \lIf{$|\SB| < |\SB^*|$}{$\SB^* \leftarrow \SB$}
}
\Return $\SB^*$\;
\caption{Monte Carlo Set Cover}
\end{algorithm}

\subsubsection*{Análise de complexidade --- Solver Randômico}

\begin{itemize}
  \item \textbf{Tempo por tentativa:} $\bigo{N_{15} \cdot N_p}$ no pior caso
        (varredura completa de $\mathcal{S}_{15}$ verificando cobertura).
        Na prática, a cobertura acumulada encurta o laço.
  \item \textbf{Tempo total ($T$ tentativas):} $\bigo{T \cdot N_{15} \cdot N_p}$.
  \item \textbf{Espaço:} $\bigtheta{N_{15} + N_p}$.
  \item \textbf{Garantia:} nenhuma de otimalidade; sem razão de aproximação provada.
        Útil como \emph{baseline} e para verificação empírica.
\end{itemize}

% ---------------------------------------------------------------------------
\section{Resultados Obtidos}
% ---------------------------------------------------------------------------

\subsection{Resultados do Greedy}

\begin{table}[H]
\centering
\caption{Resultados do algoritmo Greedy --- $n=25$, $k=15$.}
\begin{tabular}{crrrrrr}
\toprule
$p$ & $|\OPT| \geq$ (LB-LP) & $|\OPT| \geq$ (LB-Sch) & $|\SB_{\text{greedy}}|$ & Gap & Tempo (s) & $|\SB|/|\mathcal{S}_{15}|$ \\
\midrule
14 & 297.160 & 297.172 & 532.555 & 1.79$	imes$ & --- & 16.3\% \
13 & 49.527 & 58.887 & 	extit{(em execução)} & --- & --- & --- \
12 & 11.430 & 13.175 & 38.100 & 3.33$	imes$ & --- & 1.2\% \
11 & 3.266 & 3.370 & 	extit{(em execução)} & --- & --- & --- \
\bottomrule
\end{tabular}
\end{table}

\subsection{Interpretação do arquivo de resultado}

Os resultados são armazenados em arquivos \texttt{.npy} (formato NumPy binário)
contendo arrays de inteiros \texttt{uint32}. Cada inteiro é o bitmask de 25 bits
do conjunto $X \in \SB_{15,p}$. Para inspecionar:

\begin{lstlisting}[language=Python, basicstyle=\ttfamily\small,
                   frame=single, caption={Leitura e interpretação do resultado.}]
import numpy as np

# Carregar resultado de p=14
sb = np.load("resultados_SB15_14.npy")   # array uint32 com 532.555 elementos

# Converter bitmask para conjunto
def bitmask_para_conjunto(mask):
    return {e+1 for e in range(25) if mask & (1 << e)}

# Primeiro elemento de SB
print(bitmask_para_conjunto(sb[0]))
# -> {1, 2, 3, 4, 5, 6, 7, 8, 9, 10, 11, 12, 13, 14, 15}

# Verificar que Y={1..14} esta coberto por sb[0]
Y_mask = sum(1 << (e-1) for e in range(1, 15))   # {1..14}
print((sb[0] & Y_mask) == Y_mask)   # True
\end{lstlisting}

\subsection{Consistência com a teoria}

O gap observado de $1{,}79\times$ para $p=14$ satisfaz:
\[
  |\SB_{\text{greedy}}| \leq \Harm{|\mathcal{S}_{14}|} \cdot |\OPT|
  \approx 16{,}3 \cdot |\OPT|
\]
com folga confortável. O lower bound de Schönheim $|\OPT| \geq 297.172$
implica $|\OPT| / |\SB_{\text{greedy}}| \geq 0{,}558$, ou seja, a solução
greedy é no máximo $1{,}79\times$ o ótimo.

% ---------------------------------------------------------------------------
\section{Gargalos Computacionais e Escalabilidade}
% ---------------------------------------------------------------------------

\subsection{Gargalos identificados}

\begin{enumerate}
  \item \textbf{Construção dos índices hash} ($\bigtheta{N_{15} + N_p}$):
        executado uma vez; demora $\approx 8$s para $N_{15} = 3{,}27$M e $N_p = 4{,}46$M.

  \item \textbf{Loop interno de atualização de \texttt{count}}:
        $\C{15}{p} \cdot \C{25-p}{15-p}$ iterações Python por iteração principal.
        Para $p=11$: $1.366.365$ iterações Python por passo do greedy.
        Python puro executa $\approx 300{-}600$K iterações/s com dict lookups,
        resultando em $\approx 2{,}3$s por passo e $\approx 4$h no total.

  \item \textbf{Inflação do heap lazy}:
        a cada iteração, até $\C{15}{p}\cdot\C{25-p}{15-p}$ novas entradas
        são empilhadas. Para $p=13$: $6.930$ entradas/iteração
        $\times \approx 90.000$ iterações $\approx 624$M entradas acumuladas,
        degradando as operações do heap de $\bigo{\log N_{15}}$ para
        $\bigo{\log(\text{tamanho do heap})}$ maior.

  \item \textbf{Verificação de cobertura} ($\bigo{|\SB| \cdot N_p}$):
        inviável em escala completa ($532$K $\times$ $4{,}5$M $= 2{,}4$T operações).
        Solução: verificação por amostragem com broadcasting NumPy em batches
        de 512 $\times$ 10.000, concluindo em $\approx 45$s.
\end{enumerate}

\subsection{Escalabilidade}

\begin{table}[H]
\centering
\caption{Escalabilidade do Greedy por valor de $p$.}
\begin{tabular}{clll}
\toprule
$p$ & Tempo observado & Limitante teórico & Melhoria possível \\
\midrule
14 & 297.160 & 297.172 & 532.555 & 1.79$	imes$ & --- & 16.3\% \
13 & 49.527 & 58.887 & 	extit{(em execução)} & --- & --- & --- \
12 & 11.430 & 13.175 & 38.100 & 3.33$	imes$ & --- & 1.2\% \
11 & 3.266 & 3.370 & 	extit{(em execução)} & --- & --- & --- \
\bottomrule
\end{tabular}
\end{table}

\paragraph{Melhorias de escalabilidade identificadas.}
\begin{itemize}
  \item \textbf{Índice invertido pré-computado}:
        armazenar, para cada $Y \in \mathcal{S}_p$, os índices dos $X \in \mathcal{S}_{15}$
        que o contêm. Elimina o laço de extensões, mas requer
        $N_p \cdot \C{25-p}{15-p}$ entradas ($\approx 4{,}5$B para $p=11$, inviável em RAM).

  \item \textbf{Implementação em C/Cython}:
        a $\approx 300\times$ diferença de velocidade entre Python e C tornaria
        o greedy para $p=11$ executável em $\approx 50$s.

  \item \textbf{Paralelização}:
        as $\C{15}{p}$ extensões de $Y$ são independentes e paralelizáveis com
        \texttt{multiprocessing}, reduzindo o tempo proporcional ao número de núcleos.
\end{itemize}

% ---------------------------------------------------------------------------
\section{Comparação entre Estratégias}
% ---------------------------------------------------------------------------

\begin{table}[H]
\centering
\caption{Comparação entre os solvers implementados.}
\begin{tabular}{lllll}
\toprule
Solver & Complexidade tempo & Garant. solução & Escalável & Usa para \\
\midrule
Greedy & $\bigo{K \cdot \C{15}{p} \cdot \C{25-p}{15-p}}$ & $\Harm{N_p} \cdot \OPT$ & Sim & todos os $p$ \\
ILP (direto) & NP-difícil & Ótima & Não & inviável \\
ILP (pós-greedy) & Exponencial (espaço red.) & Ótima sobre cand. & Parcial & $p \in \{11,12\}$ \\
Randômico & $\bigo{T \cdot N_{15} \cdot N_p}$ & Nenhuma & Sim & baseline \\
\bottomrule
\end{tabular}
\end{table}

A estratégia adotada nos Programas 2--5 foi o Greedy, complementado pelo ILP
pós-greedy para as instâncias de menor escala ($p \in \{11, 12\}$), onde o
tamanho mínimo de $|\SB|$ (3.266 e 11.430, respectivamente) torna o número de
candidatos manejável.

% ---------------------------------------------------------------------------
\section{Conclusão}
% ---------------------------------------------------------------------------

O problema de cobertura de combinações é uma instância de \textsc{Min-Set-Cover},
provadamente NP-difícil (Teorema~\ref{thm:np}), o que justifica o uso de
heurísticas e algoritmos de aproximação.

O algoritmo Greedy (Teorema~\ref{thm:greedy-correctness}) garante cobertura
completa com razão de aproximação $\Harm{|\mathcal{S}_p|} \approx \ln|\mathcal{S}_p|$
(Teorema~\ref{thm:approx}). Para $p=14$, o gap obtido de $1{,}79\times$ está
bem abaixo da garantia teórica de $16{,}3\times$.

A análise empírica revelou que o custo por iteração cresce como
$\C{15}{p} \cdot \C{25-p}{15-p}$, explosão combinatorial que torna $p=11$
($\approx 1{,}37$M operações/iteração) o caso mais custoso apesar de exigir o
menor $|\SB|$. Esta observação não é intuitiva e emerge naturalmente da
análise de complexidade.

% ---------------------------------------------------------------------------
\bibliographystyle{plainnat}
\begin{thebibliography}{9}

\bibitem{karp1972}
Karp, R.M. (1972).
\textit{Reducibility Among Combinatorial Problems}.
In: Miller, R.E., Thatcher, J.W. (eds.), Complexity of Computer Computations.
Plenum Press, New York, pp. 85--103.

\bibitem{johnson1974}
Johnson, D.S. (1974).
Approximation algorithms for combinatorial problems.
\textit{Journal of Computer and System Sciences}, 9(3), 256--278.

\bibitem{lovasz1975}
Lovász, L. (1975).
On the ratio of optimal integral and fractional covers.
\textit{Discrete Mathematics}, 13(4), 383--390.

\bibitem{schoenheim1964}
Schönheim, J. (1964).
On coverings.
\textit{Pacific Journal of Mathematics}, 14(4), 1405--1411.

\bibitem{feige1998}
Feige, U. (1998).
A threshold of $\ln n$ for approximating set cover.
\textit{Journal of the ACM}, 45(4), 634--652.

\bibitem{dinur2014}
Dinur, I., Steurer, D. (2014).
Analytical approach to parallel repetition.
\textit{Proceedings of STOC 2014}, 624--633.

\end{thebibliography}

\end{document}
